
la série \(1+\ldots+\bar a_n+\ldots\), où \(\bar a_n\) est la classe de \(a_n\) dans \(\operatorname{Gr}^nA\). Son
noyau est le groupe \(\Ghat'(A)\) des séries formelles telles que~:
\[
  a_0=1;\qquad \forall i\geq1,\quad a_i\in\Fil^{i+1}A;
\]
on a donc un isomorphisme, fonctoriel en \(A\)~:
\[
  \Ghat(A)/\Ghat'(A)\xrightarrow{\sim}1+\widehat{\operatorname{Gr}A}^{+}.
\]

\textbf{Lemme 1.12.} \uline{Soit} \(H:\mathcal G\to\underline{\mathrm{Ab}}\,\underline{\mathrm{fil}}\) \uline{un foncteur sur} \(\mathcal G\)~; \uline{alors l'homomorphisme canonique}~:
\[
  \rho:\Hom_{\underline{\mathrm{Ab}}\,\underline{\mathrm{fil}}}(1+\widehat{\,\mathord{.}\,}^{+},H)
  \longrightarrow
  \Hom_{\underline{\mathrm{Ab}}\,\underline{\mathrm{fil}}}(\Ghat,H\circ\operatorname{Gr}),
\]
\uline{est un isomorphisme}.

L'homomorphisme \(\rho\) est défini comme étant le composé~:
\[
  \Hom(1+\widehat{\,\mathord{.}\,}^{+},H)
  \longrightarrow
  \Hom(1+\widehat{\operatorname{Gr}\,\mathord{.}\,}^{+},H\circ\operatorname{Gr})
  \longrightarrow
  \Hom(\Ghat,H\circ\operatorname{Gr}),
\]
la dernière flèche étant définie par l'homomorphisme \(\Ghat\longrightarrow1+\widehat{\operatorname{Gr}\,\mathord{.}\,}^{+}\)

On va construire un morphisme en sens inverse~:
\[
  \rho':\Hom(\Ghat,H\circ\operatorname{Gr})
  \longrightarrow
  \Hom(1+\widehat{\,\mathord{.}\,}^{+},H).
\]
Pour cela, on va montrer que tout morphisme de foncteurs de \(\Ghat\) dans \(H\circ\operatorname{Gr}\) se factorise par le morphisme canonique \(\Ghat\longrightarrow1+\widehat{\operatorname{Gr}\,\mathord{.}\,}^{+}\), ce qui équivaut à montrer qu'il s'annule sur \(\Ghat'\) d'après 1.11.

Soient \(A\in\operatorname{Ob}(\mathcal C)\), \(\chi\in\Ghat'(A)\)~; on définit un homomorphisme de
